\documentclass{article}
\usepackage[utf8]{inputenc}
\usepackage{a4wide}

\begin{document}

\section*{ช่องว่างห่างสุด}

กำหนดชุดของเลขจำนวนเต็มให้ชุดหนึ่ง

ให้หาค่าความแตกต่าง\textbf{สูงสุด}ของเลขสองตัวที่อยู่ติดกันหลังจากการเรียงลำดับแล้ว

หากเลขชุดนี้มีแค่จำนวนเดียว ให้ตอบ 0\\



ควรจะเขียนอัลกอริทึมให้อยู่ในระดับ\texttt{O(N)}



\hfill\break



\textbf{ตัวอย่าง}



\begin{quote}

\textbf{Input:}



4



19 1 12 8



\textbf{Output:}7

\end{quote}



\texttt{{[}\ 1,\ 8,\ 12,\ 19\ {]}\ //\ 8\ -\ 1\ =\ 7}



{}



\subsection*{Input}
บรรทัดแรกจำนวนขนาดของชุดตัวลข



บรรทัดที่สองตัวเลขคั่นด้วยช่องว่าง



\subsection*{Output}
{ค่าความแตกต่าง}{สูงสุด}{ของเลขสองตัวที่อยู่ติดกันหลังจากการเรียงลำดับแล้ว}\\



\end{document}
